\section{Experimental Setup}

Per garantire la replicabilità del lavoro, l'addestramento è stato meticolosamente calibrato seguendo le best practice del deep learning generativo.

\subsection{Ribilanciamento dei Set di Valutazione}

Un limite critico riscontrato nel dataset originale riguardava la partizione di validazione ufficiale, composta da soli 16 campioni. Tale dimensione esigua rendeva le metriche di validazione statisticamente insignificanti e soggette a fluttuazioni estreme durante l'addestramento. 

Per ovviare a questo problema, è stata effettuata una procedura di \textit{stratified re-sampling}. Il validation set originale è stato aumentato prendendo il 12.5\% del test set. Questa ridistribuzione ha garantito che sia il set di validazione che quello di test mantenessero la proporzione originale delle classi, fornendo al contempo una base statistica sufficientemente ampia per monitorare accuratamente la convergenza del modello e prevenire l'overfitting.

\subsection{Dinamiche di Ottimizzazione e Parametri della GAN}

La C-DCWGAN-GP opera su immagini standardizzate a una risoluzione di 128x128 pixel. L'addestramento, esteso per un arco di 100 epoche, segue una dinamica asincrona: il Critic viene aggiornato con una frequenza di 5 a 1 rispetto al Generatore $(n_{critic} = 5)$

Questo squilibrio controllato assicura che il Generatore riceva sempre gradienti densi e ben definiti da un "avversario" vicino all'ottimalità.

L'ottimizzatore Adam è stato configurato con un learning rate iniziale di $10^{-4}$ e coefficienti di momentum impostati a $\beta_1 = 0.0$ e $\beta_2 = 0.9$. L'azzeramento di $\beta$ è una scelta deliberata per eliminare l'inerzia dei gradienti storici, permettendo al modello di reagire istantaneamente ai rapidi cambiamenti del Critic. La convergenza finale è assistita da un MultiStepLR scheduler, che riduce il learning 
rate del 80\% (fattore $\gamma = 0.2$) alle epoche 60 e 80, permettendo una sintonizzazione fine delle feature microscopiche nelle radiografie.

\subsection{Training del Classificatore}

I classificatori ResNet (sia quello baseline che quello finale) sono stati addestrati con una strategia più conservativa: 10 epoche con una batch size di 32 e un learning rate fisso di 0.001. Questa configurazione è stata scelta per isolare l'effetto della qualità dei dati (l'augmentation) rispetto alla complessità dell'ottimizzazione, mantenendo un ambiente di test controllato e confrontabile.